
Creating accurate models describing how a protein carries its function, holds promise to both enable and speed-up
scientific discovery and translational research.
Since almost all protein functions derive from the inter-atomic interactions among the proteins' constituent amino-acids,
first-principles methods rooted in physics are established approaches for building models that generalize across
diverse protein functions.
Despite their promise of generalizability, physics-based methods are still greatly limited in usability
by slow speeds and high need for computational resources.

In recent years, machine learning (ML) - driven by highly expressive neural networks - has risen to fill
this gap by substituting first-principles, general modeling with effective theories learned from data,
each specific to a single or a circumscribed set of tasks and to a particular representation choice.
These models are fast, but their speed is traded for potential loss of relevant information in the
effective theory learned from data.
Thus, the main challenge in developing a useful machine learning model rests in carefully balancing
model expressivity, training data, and choice of representation, such that the learned "shortcuts" are predictive
of real-world phenomena.

In this thesis, we attempt to achieve this balance by applying the very first-principles thinking
that ML had seemed to have rendered obsolete.
Indeed, the principle that a protein successfully carries its function owing to a circumscribed set of
local inter-atomic effects, can help machine learning research just as much as molecular dynamics.
Therefore, I developed data-driven models of \textit{local atomic environments},
reasoning that the resulting effective theories would capture the very local inter-atomic effects that we know drive
protein function, thus rendering the models predictive across diverse functional settings.

Just like a fundamental principle underlying protein function is inter-atomic interactions, one of the core elements
of said interactions is their symmetry with respect to rotations and translations in 3D: the global orientation of the
protein or protein-protein complex has no bearing on the occurring interactions.
The field of Geometric Deep Learning offers a principled framework to constrain neural networks to make predictions
that respect desired symmetries without the use of additional training data.
This ensures that the effective theories learned by the models possess these symmetries, deleting at the root the risk
that they have learned symmetry-breaking shortcuts.

In this thesis, these first-principle considerations are combined into a unified neural network framework
for building \textit{rotationally-equivariant models of residue-centered atomic neighborhoods}, which is
described in detail in Chapter 2.
I then applied this framework to develop several models spanning diverse aspects of protein structure and function,
which all achieve predictive performances closely matching or exceeding the state-of-the-art,
evidencing the efficacy of the proposed framework.

In Chapter 3, I applied the SO(3)-equivariant framework to the domain of unsupervised learning, developing the Holographic (Variational) AutoEncoder (H-(V)AE).
H-(V)AE's key idea is to disentangle the latent space into a maximally informative invariant embedding - which learns a compact representation
of the input in a "canonical" orientation - and a single rotation matrix that learns the orientation of the input relative to the "canonical" one.
Through simple experiments on spherical images, we provide evidence that the "canonical" orientation is one where all the training samples are
maximally overlapped on top of each other.
Finally, when applied to protein neighborhoods H-(V)AE was able to learn compact embeddings that efficiently encode key biophysical and structural
properties of the neighborhoods, and that can be successfully used as features for downstream tasks; indeed, pairing them with a Random Forest Regressor~\citep{breiman_random_2001}
resulted in state-of-the-art performance on predicting protein-ligand binding affinity on the ATOM3D Ligand Binding Affinity dataset~\citep{townshend_atom3d_2021}.
Furthermore, H-(V)AE neighborhood embeddings are currently being used by a colleague to easily identify structural features that are indicative
of a protein's functional family, following similar principles as~\citep{sael_improved_2010, daberdaku_exploring_2018}.

Then, in Chapter 4, I used the SO(3)-equivariant framework to tackle the task known as "rotamer packing", which concerns finding physically correct side-chain
conformations that are consistent with a backbone structure and an amino-acid sequence that is predicted to fold into the backbone.
This task is common in protein design pipelines, as it complements models that design backbone structures~\citep{watson_novo_2023, lewis_scalable_2025}
and amino-acid sequences~\citep{dauparas_robust_2022} without modeling exact atomic placement of the side-chains.
The method I proposed first predicts each side-chain conformation using a local model of backbone-atom and spatially localized amino-acid labels,
and then iteratively refines the predictions using a model of the previously-predicted all-atom conformations.
The resulting algorithm - H-Packer - is not state-of-the-art on common benchmarks, but nonetheless its sufficiently high performance, speed and ease-of-use
have enabled it to be widely used by the community, especially through the popular BioEmu codebase~\citep{lewis_scalable_2025}.

In Chapter 5, I revisited the work by my former colleague and mentor, Michael Pun, on H-CNN, a model trained to predict the likelihood
of finding each amino-acid in a local atomic context. In his work, Mike provided proof-of-principle results that the learned likelihoods
could be used to zero-shot the effects of mutations on a protein's stability and on protein-protein binding affinity.
In this Chapter, I contributed several methodological improvements, culminating in the HERMES family of models.
These improvements include: a $\times 2.75$ faster and easier to optimize neural network architecture;
improved modeling of local relaxation as a result of mutations at no inference-time cost, with the model known as HERMES-\textit{amortized};
an efficient pipeline for fine-tuning HERMES on desired mutation effects, used to create HERMES-\textit{amortized}
as well as state-of-the-art models of stability effects when fine-tuning on experimental stability measurements;
a clean code repository, easy to install and use even for practitioners.
I also probed HERMES and competing models for their strengths, weaknesses,
and biases, tying them to methodological design choices. These analyses, other than serving as blueprints for the community, 
reveal that HERMES' full-atom input endows it with a strong preference for steric-preserving substitutions,
and also that pre-training on wild-type amino-acid classification introduces a wild-type preference bias that is stronger for proteins
with closer homologs in the pre-training data.
Finally, I provided evidence that HERMES-\textit{amortized} can identify mutations known to stabilize viral antigens in a conformation
amenable for vaccine development, and presented an engineering pipeline that leverages HERMES-\textit{amortized} to speed-up
vaccine development.

Finally, in Chapter 6, I showed further evidence for the impact that HERMES can have on biological discovery.
Specifically, I studied the interaction between a given T-Cell Receptor (TCR) and a peptide antigen presented
on an MHC molecule. These interactions are key events that regulate adaptive immunity, and their modeling has been largely elusive,
owing to a lack of exhaustive data, coupled with the high level of degeneracy characteristic of the system (a single TCR can interact
with many peptides, and vice-versa).
In this chapter, I used HERMES to model the distribution of peptides that are presented by a specific MHC, and recognized by a specific TCR,
a problem known among immunologists as that of TCR specificity.
Our results show that HERMES, without almost any specific knowledge~\footnote{barring what was in the pre-training data, which deduplication
at 30\% similarity greatly reduced, see Chapter~\ref{sec:hermes}}
of TCR-pMHC interactions, but access to a reliable wild-type TCR-pMHC structure, can be used to design a broad range of peptides that a target TCR
interacts with. HERMES not only successfully identifies peptide residues that do not contribute to TCR binding and thus can be ignored,
but is also able to suggest peptides that engage in stronger interactions with the TCR than the wild-type.
This chapter is not only a testament to HERMES' ability to model amino-acid substitutions in diverse structural contexts,
it is also a testament to the importance of paradigm shifts.
Indeed, prior to this work, most research on TCR specificity prediction and design was performed using sequence representations
and relied on experimental training-data distributions heavily skewed towards high TCR diversity and low peptide diversity.
Instead, using a structural representation relieves us from relying on experimental data for training, and enables
the use of a zero-shot model like HERMES. This of course comes at the expense of having a reliable TCR-pMHC structure in the first place,
which is not always available and ignores any possible alternative conformation.


% Easy-to-use implementations of these tools are publicly-available and have received increasing adoption
% from the protein engineering community at large. (this is repeating something said for H-packer already...)


The successes of the methods presented in this thesis are a testament to the effectiveness of the design choices
that underlie the structure-based, residue-centered framework used to build all of them.
This effectiveness likely hinges on two main factors.

First, the neural networks built with this framework are particularly efficient at modeling the atomic environments
to the degree that is required by the tasks considered in this thesis.
Indeed, the SO(3)-equivariant neural networks built in this thesis are strictly less expressive than several SE(3)-equivariant
alternatives commonly referenced in the atomic modeling literature~\citep{batzner_e3-equivariant_2022, liao_equiformer_2022, thomas_tensor_2018}.
Their expressiveness - and associated computational complexity - is however unnecessary for the tasks considered here.
In fact, those neural networks are often used to fit force measurements \textit{on each atom}, and thus require representing the neighborhood
and all the interactions with high atom-level precision.
By contrast, the neural networks in this thesis encode atomic environments into a lossy representation, and processes it with way fewer
expensive tensor product operations.
However, they are also only required to make predictions \textit{for the focal residue}, a generally more ``coarse-grained" setting
that likely requires less precise details.
Therefore, the success of the models in this thesis is indicative that less expressive models are sufficient for the tasks considered here.
Finding the right level of expressiveness for a task is important:
not only because lower expressiveness reduces the risk for overfitting~\citep{bishop_pattern_2006},
but also because it often reduces computational complexity, and the risk that something might go wrong during model development.
Indeed, strictly equivariant and expressive models~\citep{batzner_e3-equivariant_2022, liao_equiformer_2022, thomas_tensor_2018}
have been known to be both very computationally expensive and difficult to optimize, to the point that recent works
have been trying to relax the constraint of equivariance to improve optimization properties~\citep{pertigkiozoglou_recurrent_2026}.
Instead, I believe the framework presented in this thesis will stand the test of time owing to its balanced simplicity,
which keeps it fast and easy to optimize without compromising on equivariance.

Second, the focus on modeling local atomic environments directly results in the learning of effective theories over the majority of inter-atomic
and inter-residue interactions that render single protein folds and protein-protein interactions stable, which are hallmarks of protein function.
This is in contrast for example with models trained to model amino-acid distributions in functional protein sequences without worrying about structure,
which have been shown to capture meaningful correlates between amino-acid variability and functional variability only as deeply as evolution has sampled them~\citep{zhang_protein_2024}.

The design choices of locality and focusing on static atomic environments also pose obvious limits.
For example, as discussed in Chapters~\ref{sec:hermes}~and~\ref{sec:tcrpmhc}, these choices limit the accuracy of HERMES to mutation effects
that are driven by local interactions surrounding the mutated site, and in particular with minimal necessary side-chain relaxations upon mutations
(limit only partially mitigated by HERMES-\textit{amortized}).
However, what is most dangerous about these limits - and others that may be unveiled - is the ignorance of them.
Indeed, these limits simply define the domain in which the effective theory represented by the model is applicable.
Therefore, characterizing them provides a "how to" manual for when the model is applicable and what it can accomplish:
not much different from a function contract in software engineering.
Knowledge of these limits also turns the model into a hypothesis generation machine. For example, the success of HERMES-\textit{fixed}
on scoring and design of TCR-pMHC interactions could be seen as evidence that TCR-pMHC interactions are largely driven by local effects,
and that neural substitutability at TCR-pMHC interfaces is largely driven by preservation of steric properties 
- a hypothesis that is actually supported by a complementary line of evidence~\citep{pyo_data-driven_2025}.
Last, being able to make useful predictions with low computational requirements and high speed unlocks use cases of the model that
are not feasible with slower albeit potentially more accurate approaches.

Nevertheless, just like HERMES-\textit{amortized} improved the model's performance without changing its speed,
other improvements are possible within the framework.
A relatively small improvement would be distilling the predictions of each 10-model HERMES ensemble into a single model~\citep{hinton_distilling_2015},
to further increase inference speed and make the models even more CPU-friendly.
I would also want to estimate prediction uncertainty for both HERMES and H-Packer. This would have the primary purpose of guiding users'
expectations with respect to model predictions, but also the predicted uncertainty itself may be predictive of relevant physical quantities,
for example side-chain conformational entropy in the case of H-Packer.
A natural place to start would be reporting and analyzing the 10-model ensemble variance of HERMES.\\
Thinking instead about large-scale extensions of the framework, insightful future work would be training HERMES without the full side-chains,
but still with local neighborhoods, and using an input parameterization similar to that of the \textit{initial guess} H-Packer model.
This would disentangle locality and all-atom design choices, resulting in a model with inductive biases "in-between" those of HERMES
and of ProteinMPNN.
Finally, I envision the efforts in this thesis to potentially be unified in a single model that performs auto-regressive prediction of both
amino-acid labels and side-chain configurations, with the additional inclusion in the neighborhood featurization of atomic types beyond those
found in proteins - like phosphorus and common metallic ions.
Such a model would still only use local atomic neighborhoods as input, but it would include "mask" and amino-acid-specific tokens to handle
the cases in which amino-acid labels and side-chain conformations are respectively unknown; in many ways it would be a direct competitor of
LigandMPNN~\citep{dauparas_atomic_2025}, but developed with strict locality and our rotationally equivariant framework.
This model would be a direct substitute of both HERMES and H-Packer, and in addition be able to
(1) design structure-conditioned sequences without assuming side-specific independence,
and (2) better design protein interfaces against non-protein molecules like DNA, ions, and small-molecule ligands.
